\documentclass[11pt]{article}

\usepackage[margin=1in]{geometry}
\usepackage{amsmath,amssymb}
\usepackage{enumitem}
\usepackage[most]{tcolorbox}
\usepackage{xcolor}

\setlength{\parindent}{0pt}
\setlength{\parskip}{0.6em}


% BOX STYLES


\newtcolorbox{bluebox}{
colback=blue!7,
colframe=blue!70!black,
boxrule=1pt,
arc=4pt,
breakable,
enhanced
}

\newtcolorbox{purplebox}{
colback=violet!7,
colframe=violet!70!black,
boxrule=1pt,
arc=4pt,
breakable,
enhanced
}

\newtcolorbox{redbox}{
colback=red!7,
colframe=red!70!black,
boxrule=1pt,
arc=4pt,
breakable,
enhanced
}


% DOCUMENT


\begin{document}

\title{\textbf{Algebra 1 Handout \\ \small The Big Picture}}
\author{VPSS}
\date{March 2026}

\maketitle
\hrule
\vspace{0.8em}


% LEARNING TARGETS


\section*{Learning Targets}

\begin{itemize}[leftmargin=1.4em]
\item I can identify and use key Algebra 1 vocabulary.
\item I can solve equations and justify each step.
\item I can model real-world situations with algebraic expressions or equations.
\end{itemize}


% PHILOSOPHY SECTION


\begin{bluebox}

\section{Philosophy, Meaning, Purpose, Quirks}

\textbf{What is algebra really?}

Algebra is the language of patterns and relationships. Instead of working only with specific numbers, we introduce \emph{symbols} (usually letters) that represent numbers we do not yet know or numbers that may vary.

This lets us express general truths rather than isolated examples.

For example

\[
a+b=b+a
\]

This is not about two specific numbers. It states a rule that works for \emph{all} numbers. Algebra compresses infinitely many numerical statements into a single symbolic statement.

\subsection*{Why letters?}

Letters serve several purposes in mathematics.

\textbf{Unknown quantities}

They can represent a number we want to find.

\[
x+5=12
\]

\textbf{Variables}

They can represent quantities that change.

\[
d = rt
\]

Distance depends on rate and time.

\textbf{General mathematical rules}

Letters allow us to state universal identities.

\[
(a+b)^2 = a^2 + 2ab + b^2
\]

Without symbols, mathematics would be limited to individual calculations instead of general reasoning.

\subsection*{The philosophy of contest math}

Contest mathematics treats algebra as a \textbf{tool for thinking}, not just calculating.

In many problems, the key step is choosing the right representation. Difficult problems often become simple once the situation is translated into the right algebraic expression.

Strong algebraic thinking involves

\begin{itemize}
\item translating word problems into equations
\item introducing useful variables
\item recognizing patterns and symmetry
\item simplifying expressions strategically
\end{itemize}

\subsection*{Meaning over manipulation}

Algebra is not just a list of mechanical rules. Every step in solving an equation represents a logical operation.

Example:

\[
2x+6=14
\]

Subtract $6$ from both sides:

\[
2x=8
\]

Divide both sides by $2$:

\[
x=4
\]

Each operation preserves the equality.

\subsection*{Quirks and conventions}

Algebra uses a few conventions that may seem unusual at first.

\begin{itemize}

\item \textbf{Multiplication by adjacency}

\[
3x = 3\cdot x
\]

\item \textbf{Variable names do not matter}

\[
x+2=5 \quad \text{means the same as} \quad y+2=5
\]

\item \textbf{Order properties}

\[
a+b=b+a
\]

but

\[
a-b \ne b-a
\]

\item \textbf{Expressions vs equations}

\[
x^2+3x
\]

is an expression.

\[
x^2+3x=10
\]

is an equation.

\end{itemize}

\subsection*{Why this matters for contest math}

Strong problem solvers use algebra as a creative language.

Instead of asking

\begin{center}
“What formula should I apply?”
\end{center}

they ask

\begin{itemize}

\item What quantities should I represent with variables?
\item What relationships exist between them?
\item Can this expression be rewritten more simply?

\end{itemize}

The best algebraic move is often the one that reveals hidden structure.

\[
\text{Algebra = expressing ideas precisely and manipulating them intelligently}
\]

\end{bluebox}


% EXAMPLE


\section*{Example for LaTeX and General Problems}

\textbf{Solve}

\[
3(2x-1)=15
\]

\textbf{Solution}

\[
\begin{aligned}
3(2x-1) &= 15 \\
6x-3 &= 15 \\
6x &= 18 \\
x &= 3
\end{aligned}
\]

\textbf{Check}

\[
3(2\cdot3-1)=3(5)=15 \quad \checkmark
\]

\begin{purplebox}
\section{General Tips for Learning}

\begin{itemize}
\item Always understand \textbf{why} a step works, not just how to do it.
\item Try to rewrite problems in your own words before solving.
\item Introduce variables carefully and define what they represent.
\item Check answers by substituting them back into the original equation.
\item Look for patterns and structure rather than expanding everything immediately.
\end{itemize}

\section{Lesson 1: The Basic Tools for Simple Problems}

\textbf{Objective:} Single-variable algebra can be understood intuitively.

In many algebra problems, we are solving for a \textbf{single unknown quantity}.  
The goal is to isolate that variable using operations that preserve equality.

Key tools include:

\begin{itemize}

\item \textbf{Balancing equations}

\[
x+5=12
\]

Subtract $5$ from both sides:

\[
x=7
\]

\item \textbf{Distributing}

\[
3(x+2)=3x+6
\]

\item \textbf{Combining like terms}

\[
2x+3x=5x
\]

\item \textbf{Inverse operations}
\[
Addition \iff Subtraction \\ 
Multiplication \iff Division
\]
\end{itemize}

Most introductory algebra problems can be solved using only these few tools.

\end{purplebox}

% PROBLEMS

\begin{redbox}

\section{Problems}

Complete the following independently. Show clear work for full credit.

\begin{enumerate}[leftmargin=1.5em,label=\textbf{\arabic*.}]

\item $9x-4=41$

\item $7(x+2)=3x+30$

\item $2(3x-5)-4=18$

\item $\dfrac{5x-1}{2}=12$

\item A phone plan charges \$15 base fee plus \$8 per gigabyte.  
Write and solve an equation for the number of gigabytes used if the total bill is \$71.

\item The sum of two consecutive integers is $47$.  
Write and solve an equation to find the integers.

\end{enumerate}

\end{redbox}
% EXIT TICKET
\section*{Exit Ticket}

In \textbf{2--3 sentences}: What step do you find most important when solving equations, and why?

\vspace{0.8in}
\hrule \\*
\textbf{Part of a 3-Piece Handout Series on Algebra I} \\
\textbf{If you plan to use this, cite me!} \\
\textbf{Reach out at cowboyhehe1@gmail.com}
\end{document}